\documentclass[12pt,a4paper]{report}


\usepackage[utf8]{inputenc}
\usepackage[T1]{fontenc}
\usepackage[french]{babel}
\usepackage{geometry}
\geometry{margin=2.5cm}

\usepackage{graphicx} 

\begin{document}
	
	
	\begin{titlepage}
		\centering
		
	
		\includegraphics[width=0.25\textwidth]{logo.png}\\[2cm]
		
	
		{\Huge \textbf{Théories et Pratiques de
				l’Investigation Numérique}}\\[1.5cm]
		
		
		{\Large Matière : Théories et Pratiques de
			l’Investigation Numérique}\\[0.5cm]
		{\Large Enseignant :M. MINKA Thierry}\\[0.5cm]
		{\Large Année académique : 2025--2026}\\[2cm]
		
		
		{\Large  Par l'étudiante : ZE MVONDO Océanne M.}\\[0.5cm]
		{\Large Filière : 4e année CIN }\\[2cm]
		
		
		
	 \includegraphics[width=0.6\textwidth]{C:/Users/USER/Documents/COURS/CIN 4/forensic.jpg}\\[4cm]
		
	
	
	\end{titlepage}
	
	
	
	\setlength{\parskip}{1em} 
	\setlength{\parindent}{0pt} 
	

	
 
	
\section*{Expertise 2 : Ordonnance de renvoi
	}

\section*{Analyse des hypothèses sur le meurtre de la victime}

Sur la base des éléments d’investigation numérique transmis dans l’ordonnance de renvoi, à savoir plusieurs géolocalisations des suspects et de la victime, quelques conversations WhatsApp ainsi que certains listings d’appels, nous avons formulé trois hypothèses principales pour expliquer le meurtre de la victime, journaliste, dans un contexte où les suspects sont des opérateurs économiques, hommes politiques et hautes personnalités.

\begin{hypothesis}
	Le meurtre pourrait être directement lié à une enquête sensible menée par la victime, visant à révéler des pratiques illégales, des affaires de corruption ou des manipulations politiques impliquant certains opérateurs économiques ou hommes politiques. En tant que journaliste, la victime aurait pu obtenir ou être sur le point de publier des informations compromettantes susceptibles de nuire à des intérêts puissants. Ce contexte pourrait expliquer une volonté d’éliminer la victime pour empêcher la diffusion de ces révélations et protéger les acteurs impliqués.
\end{hypothesis}

\begin{hypothesis}
	Une autre hypothèse envisage que le meurtre soit le résultat d’un conflit d’intérêts entre différents groupes d’influence cherchant à contrôler l’opinion publique ou à orienter le discours médiatique. La victime, en sa qualité de journaliste, aurait pu jouer un rôle central dans la médiation ou la diffusion d’informations sensibles, devenant ainsi un obstacle à certains projets politiques ou économiques d’envergure. Le crime pourrait alors s’inscrire dans une stratégie visant à neutraliser un acteur clé capable de perturber des alliances ou des stratégies de pouvoir.
\end{hypothesis}

\begin{hypothesis}
	Enfin, le meurtre pourrait être le point d’aboutissement d’une opération délibérée de déstabilisation orchestrée par des hautes personnalités, dans le but d’envoyer un message fort et intimidant à la presse indépendante et aux voix critiques du pouvoir en place. Dans ce scénario, l’assassinat de la victime ne serait pas seulement un acte isolé, mais une démonstration de force destinée à instaurer un climat de peur et à restreindre la liberté d’expression, particulièrement dans un contexte politique tendu où la liberté de la presse est menacée.

	
	
\end{document}